\begin{description}
\item[(a)] The plate scale is simply the angular size of the image per unit
length as projected at the focal plane. Hence, for small angles,
\[ \tan\theta_i = \frac{\ell_i}{f} \approx \theta_i \]
\hfill(1.0 points)\\
where $\theta_i$ is the angular size of the image, $\ell_i$ is the unit
length, and $f$ is the focal length of the telescope, which is given by
($D \times$ F-ratio). \hfill(0.5 points)
\begin{gather*}
\theta_i = \frac{1\,\mathrm{mm}}{5 \times 305\,\mathrm{mm}}
  = 6.56 \times 10^{-4}\,\mathrm{rad} = 2.254' \\
PS \approx 2.25\,\mathrm{arcmin/mm}
\end{gather*}
\hfill(1.5 points)

\item[(b)] Given the plate scale, the diameter of the image on the focal
plane is
\[ d_{GC} = \frac{72.0''}{2.25' \times 60} = 0.532\,\mathrm{mm} \]
\hfill(1.0 points)

This implies that the area covered by the image is
\[ S_{GC} = \frac{\pi}{4} \cdot d_{GC}^2 = 0.223\,\mathrm{mm}^2 \]
which can be divided by the area of a single pixel to estimate the pixel
coverage
\[ n_p = \frac{0.223}{(3.8 \times 10^{-3})^2}
   \approx 15\,400~\text{pixels} \]
\hfill(2.0 points)

Since there are other methods to estimate such quantity, the accepted range
is
\[ 15300 \leqslant n_p \leqslant 15600 \]
\hfill(1.0 points)

\item[(c)] The signal-to-noise ratio is given by the expression
\begin{align*}
S/N &= \frac{N_{GC}}
  {\sqrt{\sigma_{GC}^2 + \sigma_{sky}^2 + \sigma_{RON}^2 + \sigma_{DN}^2}} \\
  &= \frac{N_{GC}}
  {\sqrt{N_{GC} + N_{sky} + n_p \cdot RON^2 + n_p \cdot DN \cdot t}}
\end{align*}
\hfill(3.0 points)\\
where,
\[
\begin{cases}
N_{GC} \rightarrow~\text{Total Source Count} \\
\sigma_{GC} = \sqrt{N_{GC}} \rightarrow~\text{Poisson Noise} \\
N_{sky} \rightarrow~\text{Total Sky Count} \\
\sigma_{sky} = \sqrt{N_{sky}} \rightarrow~\text{Sky Noise} \\
\sigma_{RON} = \sqrt{n_p \cdot 1 \cdot RON^2}
  \rightarrow~\text{Readout Noise} \\
\sigma_{DN} = \sqrt{n_p \cdot DN \cdot t} \rightarrow~\text{Dark Noise}
\end{cases}
\]

The noise values associated with the CCD operation can be easily calculated:
\[
\begin{cases}
\sigma_{RON}^2 = 15400 \times 5^2 = 385000~\text{counts} \\
\sigma_{DN}^2 = 15400 \times 6 \times \dfrac{1}{60} \times 15
  = 23100~\text{counts}
\end{cases}
\]
\hfill(2.0 points)

Now, the apparent visual magnitude of the globular cluster must be
calculated,
\[ V_{GC} - m_V = -2.5\log\left(\frac{\Omega_{GC}}{1}\right) \]
with $\Omega_{GC}$ being the solid angle subtended by the globular cluster,
\[ \Omega_{GC} = \pi \cdot \left(\frac{\theta}{2}\right)^2
   = \pi \times \left(\frac{72.0''}{2}\right)^2
   = 4071.5\,\mathrm{arcsec}^2 \]
\hfill(1.0 points)

Hence,
\[ V_{GC} = 20.6 - 2.5\log(4071.5) = 11.6\,\mathrm{mag} \]
\hfill(1.0 points)\\
and the associated flux of photons can now be evaluated using the
0-magnitude as a reference:
\begin{align*}
V_{GC} - V_0 &= -2.5\log\left(\frac{\phi_{GC}}{\phi_0}\right) \\
\phi_{GC} &= \phi_0 \cdot 10^{-V_{GC}/2.5} \\
  &\approx 0.234\,\mathrm{counts/nm/cm^2/s}
\end{align*}
\hfill(1.0 points)

With the calculated parameters, the total source count is:
\begin{align*}
N_{GC} &= \eta \cdot \phi_{GC} \cdot \frac{\pi}{4} \cdot D^2
  \cdot \Delta\lambda \cdot t \\
N_{GC} &= 0.80 \times 0.234 \times \frac{\pi}{4} \times (30.5)^2
  \times 88.0 \times 15 \\
  &= 180\,540~\text{counts}
\end{align*}
\hfill(1.0 points)

By referring to the signal-to-noise ratio expression and substituting the
values,
\begin{align*}
225 &= \frac{180540}{\sqrt{180540 + N_{sky} + 385000 + 23100}} \\
N_{sky} &= \left(\frac{180540}{225}\right)^2 - 385000 - 180540 - 23100 \\
N_{sky} &= 55250~\text{counts}
\end{align*}
\hfill(1.0 points)

The inverse procedure to determine the Total Source Count is taken:
\begin{align*}
\phi_{sky} &= \frac{N_{sky}}
  {\eta \cdot \frac{\pi}{4} \cdot D^2 \cdot \Delta\lambda \cdot t} \\
  &= \frac{55250}
  {0.80 \times \frac{\pi}{4} \times 30.5^2 \times 88.0 \times 15} \\
\phi_{sky} &= 7.16 \times 10^{-2}\,\mathrm{counts/nm/cm^2/s}
\end{align*}
\hfill(1.0 points)\\
which implies
\[ V_{sky} = V_0 - 2.5\log\left(\frac{\phi_{sky}}{\phi_0}\right)
   = 12.9\,\mathrm{mag} \]
\hfill(1.0 points)\\
and finally, since $\Omega_{sky} = \Omega_{GC}$
\begin{align*}
m_{sky} &= V_{sky} + 2.5\log\left(\frac{\Omega_{sky}}{1}\right) \\
  &= 12.8 + 2.5\log(4071.5) \\
% sic: the source prints 12.8 here although V_sky = 12.9 mag above.
m_{sky} &= 21.9\,\mathrm{mag/arcsec^2}
\end{align*}
\hfill(1.0 points)

Considering the accepted range of pixels from the previous item, the
accepted answers for the brightness of the sky fall within the interval
\[ 21.7 \leqslant m_{sky} \leqslant 22.1\,\mathrm{mag/arcsec^2} \]
\end{description}
